\documentclass[11pt, a4paper]{article}

% ── Packages ──────────────────────────────────────────────────────────────────
\usepackage[a4paper, top=2.5cm, bottom=2.5cm, left=2.5cm, right=2.5cm, headheight=28pt]{geometry}
\usepackage[utf8]{inputenc}
\usepackage[T1]{fontenc}
\usepackage{lmodern}
\usepackage{booktabs}
\usepackage{array}
\usepackage{tabularx}
\usepackage{longtable}
\usepackage{multirow}
\usepackage{colortbl}
\usepackage{xcolor}
\usepackage{titlesec}
\usepackage{enumitem}
\usepackage{hyperref}
\usepackage{fancyhdr}
\usepackage{parskip}
\usepackage{listings}
\usepackage[most]{tcolorbox}
\usepackage{fontawesome5}
\usepackage{microtype}
\usepackage{graphicx}

% ── Colors (KMUTNB brand) ──────────────────────────────────────────────────────
\definecolor{kmutnbBlue}{RGB}{0,56,116}
\definecolor{kmutnbGold}{RGB}{186,146,36}
\definecolor{lightGray}{RGB}{245,245,245}
\definecolor{midGray}{RGB}{200,200,200}
\definecolor{codeGray}{RGB}{248,248,248}
\definecolor{codeFrame}{RGB}{220,220,220}
\definecolor{tipteal}{RGB}{0,105,92}
\definecolor{warnAmber}{RGB}{121,85,0}
\definecolor{checkGreen}{RGB}{27,94,32}

% ── Section formatting ─────────────────────────────────────────────────────────
\titleformat{\section}
  {\large\bfseries\color{kmutnbBlue}}
  {\thesection.}{0.5em}{}
  [\color{kmutnbGold}\titlerule]

\titleformat{\subsection}
  {\normalsize\bfseries\color{kmutnbBlue}}
  {\thesubsection.}{0.5em}{}

\titleformat{\subsubsection}
  {\small\bfseries\color{kmutnbBlue}}
  {\thesubsubsection.}{0.5em}{}

\titlespacing*{\section}{0pt}{1.4ex plus .2ex}{0.6ex plus .1ex}
\titlespacing*{\subsection}{0pt}{1.0ex plus .2ex}{0.4ex plus .1ex}

% ── Header / Footer ────────────────────────────────────────────────────────────
\pagestyle{fancy}
\fancyhf{}
\fancyhead[L]{\small\color{kmutnbBlue}\textbf{KMUTNB -- Faculty of Engineering}}
\fancyhead[R]{\small\color{kmutnbBlue}Dept.\ of Electrical and Computer Engineering}
\fancyfoot[L]{\small\color{kmutnbBlue}Real-Time Operating Systems -- M.Eng.\ ECE}
\fancyfoot[C]{\small\thepage}
\fancyfoot[R]{\small\color{kmutnbBlue}Lab 02}
\renewcommand{\headrulewidth}{0.4pt}
\renewcommand{\headrule}{\hbox to\headwidth{%
  \color{kmutnbGold}\leaders\hrule height\headrulewidth\hfill}}
\renewcommand{\footrulewidth}{0.4pt}
\renewcommand{\footrule}{\hbox to\headwidth{%
  \color{midGray}\leaders\hrule height\footrulewidth\hfill}}

% ── Hyperlinks ─────────────────────────────────────────────────────────────────
\hypersetup{
  colorlinks = true,
  linkcolor  = kmutnbBlue,
  urlcolor   = kmutnbBlue,
  citecolor  = kmutnbBlue,
  pdftitle   = {Lab 02 — FreeRTOS Basics: Tasks, Queues, Semaphores and Timers},
  pdfauthor  = {KMUTNB RTOS Course},
}

% ── Listings (C and bash) ──────────────────────────────────────────────────────
\lstdefinestyle{cstyle}{
  language        = C,
  basicstyle      = \ttfamily\small,
  keywordstyle    = \bfseries\color{kmutnbBlue},
  commentstyle    = \itshape\color{gray},
  stringstyle     = \color{tipteal},
  numberstyle     = \tiny\color{midGray},
  numbers         = left,
  numbersep       = 6pt,
  stepnumber      = 1,
  backgroundcolor = \color{codeGray},
  frame           = single,
  framerule       = 0.4pt,
  rulecolor       = \color{codeFrame},
  breaklines      = true,
  showstringspaces= false,
  tabsize         = 4,
  xleftmargin     = 1.5em,
  xrightmargin    = 0.5em,
  aboveskip       = 0.8em,
  belowskip       = 0.8em,
}

\lstdefinestyle{bashstyle}{
  language        = bash,
  basicstyle      = \ttfamily\small,
  keywordstyle    = \color{kmutnbBlue},
  commentstyle    = \itshape\color{gray},
  backgroundcolor = \color{black!5},
  frame           = single,
  framerule       = 0.4pt,
  rulecolor       = \color{codeFrame},
  breaklines      = true,
  showstringspaces= false,
  xleftmargin     = 1.5em,
  xrightmargin    = 0.5em,
  aboveskip       = 0.6em,
  belowskip       = 0.6em,
}

\lstdefinestyle{textstyle}{
  basicstyle      = \ttfamily\footnotesize,
  backgroundcolor = \color{black!4},
  frame           = single,
  framerule       = 0.4pt,
  rulecolor       = \color{codeFrame},
  breaklines      = true,
  xleftmargin     = 1.5em,
  xrightmargin    = 0.5em,
  aboveskip       = 0.6em,
  belowskip       = 0.6em,
}

% ── Callout boxes ──────────────────────────────────────────────────────────────
\tcbuselibrary{skins, breakable}

\newtcolorbox{tipbox}[1][]{
  enhanced, breakable,
  colback=tipteal!8, colframe=tipteal, coltitle=white,
  fonttitle=\bfseries\small, title={\faLightbulb\quad Tip#1},
  left=6pt, right=6pt, top=4pt, bottom=4pt,
  boxrule=0.8pt, arc=3pt,
}

\newtcolorbox{warnbox}[1][]{
  enhanced, breakable,
  colback=warnAmber!8, colframe=kmutnbGold, coltitle=white,
  fonttitle=\bfseries\small, title={\faExclamationTriangle\quad Note#1},
  left=6pt, right=6pt, top=4pt, bottom=4pt,
  boxrule=0.8pt, arc=3pt,
}

\newtcolorbox{checkpointbox}[1][]{
  enhanced,
  colback=kmutnbBlue!6, colframe=kmutnbBlue, coltitle=white,
  fonttitle=\bfseries\small, title={\faCheckCircle\quad Checkpoint#1},
  left=6pt, right=6pt, top=4pt, bottom=4pt,
  boxrule=0.8pt, arc=3pt,
}

\newtcolorbox{questionbox}[1][]{
  enhanced, breakable,
  colback=kmutnbGold!10, colframe=kmutnbGold, coltitle=white,
  fonttitle=\bfseries\small, title={\faQuestion\quad Question#1},
  left=6pt, right=6pt, top=4pt, bottom=4pt,
  boxrule=0.8pt, arc=3pt,
}

% ── Checkbox list ──────────────────────────────────────────────────────────────
\newlist{checklist}{itemize}{1}
\setlist[checklist]{label=\faSquare[regular], leftmargin=1.5em, noitemsep}

% ── Misc helpers ───────────────────────────────────────────────────────────────
\newcommand{\file}[1]{\texttt{\small #1}}
\newcommand{\api}[1]{\texttt{\small #1}}

% ══════════════════════════════════════════════════════════════════════════════
\begin{document}

% ── Title block ───────────────────────────────────────────────────────────────
\begin{center}
  {\color{kmutnbBlue}\rule{\linewidth}{2pt}}\\[6pt]
  {\LARGE\bfseries\color{kmutnbBlue} Laboratory 02}\\[4pt]
  {\large\bfseries FreeRTOS Basics: Tasks, Queues, Semaphores \& Timers}\\[2pt]
  {\normalsize Real-Time Operating Systems \enspace\textbullet\enspace
   M.Eng.\ ECE \enspace\textbullet\enspace KMUTNB}\\[8pt]
  {\color{kmutnbGold}\rule{\linewidth}{1pt}}
\end{center}

\vspace{0.3cm}

\begin{tabular}{@{}p{3.8cm} p{10cm}@{}}
  \textbf{Platform}     & NXP FRDM-MCXN236 (ARM Cortex-M33 @ 96\,MHz) \\[2pt]
  \textbf{RTOS}         & FreeRTOS (bare-metal, no SDK) \\[2pt]
  \textbf{Toolchain}    & ARM GNU Toolchain (\texttt{arm-none-eabi-gcc}) + Make \\[2pt]
  \textbf{Tools}        & pyocd, serial terminal (115200\,8N1) \\[2pt]
  \textbf{Estimated time} & 2--3 hours \\[2pt]
  \textbf{CLO mapping}  & CLO\,3 (FreeRTOS implementation) \enspace CLO\,4 (synchronisation) \\
\end{tabular}

\vspace{0.4cm}

% ── Objectives ────────────────────────────────────────────────────────────────
\section{Objectives}

By completing this lab you will be able to:

\begin{enumerate}[noitemsep, leftmargin=*]
  \item Build a multi-task FreeRTOS application that uses four kernel primitives ---
        task, queue, semaphore, and software timer.
  \item Explain the \textbf{producer--consumer} pattern and observe what happens when
        a queue overflows.
  \item Use a \textbf{mutex} to protect a shared resource and demonstrate the
        corruption that occurs without one.
  \item Distinguish a \textbf{binary} semaphore from a \textbf{counting} semaphore by
        observing lost signals.
  \item Configure a \textbf{software timer} (one-shot vs.\ auto-reload) and explain the
        rules for timer callbacks.
  \item Contrast \api{vTaskDelay} and \api{vTaskDelayUntil} and measure timing drift.
\end{enumerate}

% ── Prerequisites ─────────────────────────────────────────────────────────────
\section{Prerequisites}

You must have completed \textbf{Lab 01} before starting:

\begin{checklist}
  \item ARM GNU Toolchain installed and working (\texttt{make} succeeds in
        \file{lab01\_setup\_verify/})
  \item \texttt{pyocd} installed (\texttt{pip install pyocd})
  \item FRDM-MCXN236 connected via the \textbf{MCU-LINK USB} port
  \item A serial terminal (\texttt{screen}, \texttt{minicom}, PuTTY, or Tera Term)
        at \textbf{115200\,8N1}
\end{checklist}

\begin{tipbox}
  No SEGGER SystemView is required for this lab --- all observation is done over UART.
\end{tipbox}

% ── Project Structure ─────────────────────────────────────────────────────────
\section{Project Structure}

\begin{lstlisting}[style=textstyle]
lab02_freertos_basics/
  Makefile                  <- build, flash, and debug targets
  linker_MCXN236.ld         <- flash/RAM memory map
  src/
    startup_MCXN236.s       <- vector table + Reset_Handler
    main.c                  <- hw_init, RTOS objects, scheduler start
    lab_tasks.c / .h        <- the three tasks + timer callback
    FreeRTOSConfig.h        <- kernel configuration
    uart.h / uart.c         <- bare-metal LPUART4 driver
  FreeRTOS-Kernel/          <- (cloned by `make setup`)
\end{lstlisting}

\subsection{What the code does --- the ``Sensor Data Pipeline''}

This lab simulates a small sensor pipeline.  Three tasks and one software timer
cooperate through shared kernel objects:

\renewcommand{\arraystretch}{1.3}
\begin{tabular}{@{} l c p{3.4cm} p{5.6cm} @{}}
  \toprule
  \textbf{Task / Object} & \textbf{Priority} & \textbf{Trigger} & \textbf{What it does} \\
  \midrule
  \texttt{vSensorTask}  & 3 (HIGH) & 500\,ms period (\api{vTaskDelayUntil}) &
    Generates a fake ADC reading, sends it to \api{xSensorQueue} \\
  \texttt{vDisplayTask} & 2 (MED)  & Blocks on \api{xSensorQueue} &
    Receives a reading and prints it (under \api{xUartMutex}) \\
  \texttt{vStatusTask}  & 1 (LOW)  & Blocks on \api{xStatusSem} &
    Prints uptime + queue depth (under \api{xUartMutex}) \\
  \texttt{StatusTimer}  & --- (timer task) & 1000\,ms, auto-reload &
    Callback gives \api{xStatusSem} to wake \texttt{vStatusTask} \\
  \bottomrule
\end{tabular}
\renewcommand{\arraystretch}{1}

\medskip
\textbf{Data flow:}

\begin{lstlisting}[style=textstyle]
                  xSensorQueue (depth 8)
   vSensorTask  ----------------------->  vDisplayTask
    (prio 3)         SensorReading_t        (prio 2)
    500 ms                                      |
                                                | xUartMutex
   StatusTimer ---xStatusSem--> vStatusTask -----+
    (1000 ms)                    (prio 1)        |
                                                 v
                                                UART
\end{lstlisting}

\begin{itemize}[noitemsep, leftmargin=*]
  \item \textbf{Queue} (\api{xSensorQueue}) decouples the producer (\texttt{Sensor})
        from the consumer (\texttt{Display}).
  \item \textbf{Semaphore} (\api{xStatusSem}) is a pure \emph{signal} --- the timer
        signals, the task waits.
  \item \textbf{Mutex} (\api{xUartMutex}) guarantees only one task drives the UART at
        a time.
\end{itemize}

% ── Part A ─────────────────────────────────────────────────────────────────────
\section{Part A --- Build and Flash}

\subsection{Step 1 --- Clone the FreeRTOS kernel}

\begin{lstlisting}[style=bashstyle]
cd labs/lab02_freertos_basics
make setup
\end{lstlisting}

This runs \texttt{git clone --depth=1} into \file{FreeRTOS-Kernel/}.  You only do this
once.

\begin{tipbox}
  Unlike Lab 01, this lab's Makefile also compiles \file{timers.c} --- the
  software-timer module --- because we use a software timer.
\end{tipbox}

\subsection{Step 2 --- Build}

\begin{lstlisting}[style=bashstyle]
make
\end{lstlisting}

Expected output (sizes will vary slightly):

\begin{lstlisting}[style=textstyle]
arm-none-eabi-gcc  ...  -o lab02.elf
arm-none-eabi-objcopy -O binary lab02.elf lab02.bin
   text    data     bss     dec     hex filename
  1xxxx     xxx   xxxxx   xxxxx    xxxx lab02.elf
\end{lstlisting}

\begin{warnbox}
  If \texttt{arm-none-eabi-gcc: not found}, set the toolchain path at the top of
  \file{Makefile}:
  \begin{lstlisting}[style=bashstyle, aboveskip=0pt, belowskip=0pt]
  TOOLCHAIN ?= /path/to/arm-none-eabi/bin
  \end{lstlisting}
\end{warnbox}

\subsection{Step 3 --- Flash}

\begin{lstlisting}[style=bashstyle]
make flash
\end{lstlisting}

pyocd erases, programs, and resets the MCU.  The board runs immediately.

% ── Part B ─────────────────────────────────────────────────────────────────────
\section{Part B --- UART Output Verification}

Open a serial terminal:

\begin{lstlisting}[style=bashstyle]
ls /dev/tty.usbmodem*        # find the port
screen /dev/tty.usbmodemXXXX 115200
\end{lstlisting}

Press the board \textbf{RESET} button to see the banner from the start.

\subsection{Expected output}

\begin{lstlisting}[style=textstyle]
=== RTOS Lab 02: FreeRTOS Basics ===
Tasks: Sensor | Display | Status
Primitives: Queue | Mutex | Semaphore | Timer

[MAIN] Starting scheduler...

[DISPLAY] t=0 ms  value=7
[DISPLAY] t=500 ms  value=14

[STATUS] uptime=1000 ms  queue_waiting=0
[DISPLAY] t=1000 ms  value=21
[DISPLAY] t=1500 ms  value=28

[STATUS] uptime=2000 ms  queue_waiting=0
\end{lstlisting}

\begin{tipbox}
  \textbf{Key observations:}
  \begin{itemize}[noitemsep]
    \item Two \texttt{[DISPLAY]} lines appear for every one \texttt{[STATUS]} line ---
          500\,ms vs.\ 1000\,ms.
    \item \texttt{queue\_waiting=0} --- the consumer keeps up with the producer, so the
          queue stays empty.
    \item The \texttt{value} field steps by 7 each reading
          (\texttt{raw = (raw + 7) \% 101}).
    \item Lines never interleave --- the mutex serialises UART access.
  \end{itemize}
\end{tipbox}

\begin{checkpointbox}[~A]
  Show a terminal capture of at least 5 seconds of correct output to your instructor
  before continuing.
\end{checkpointbox}

% ── Part C ─────────────────────────────────────────────────────────────────────
\section{Part C --- Experiments}

Work through these in order.  For each experiment: make the edit, run
\texttt{make flash}, observe the UART, then answer the question.

\begin{warnbox}
  Restore the original code after each experiment unless told otherwise.
\end{warnbox}

\subsection{Experiment 1 --- The queue as a producer--consumer buffer}

The producer (\texttt{Sensor}) runs every 500\,ms.  The consumer (\texttt{Display})
normally keeps up.  Now make the consumer \textbf{slower than} the producer.

In \file{lab\_tasks.c}, inside \texttt{vDisplayTask}, add a delay after the print:

\begin{lstlisting}[style=cstyle]
xSemaphoreGive(xUartMutex);
vTaskDelay(pdMS_TO_TICKS(900));   /* <-- ADD: slow consumer */
\end{lstlisting}

Rebuild, flash, and watch the \texttt{queue\_waiting} count in the \texttt{[STATUS]}
lines.

\begin{questionbox}[~1]
  \begin{enumerate}[label=(\alph*), noitemsep, leftmargin=*]
    \item The producer makes 2 readings/s; the consumer now drains $\sim$1.1/s.
          Watch \texttt{queue\_waiting} climb.  The queue depth is \textbf{8}
          (see \api{xQueueCreate} in \file{main.c}) --- roughly how many seconds until
          you see \texttt{[SENSOR] Queue full --- reading dropped!}?
    \item \api{xQueueSend} in \texttt{vSensorTask} uses a timeout of \texttt{0}.
          What would change if it used \api{portMAX\_DELAY} instead?  Would readings
          still be dropped?
    \item In one sentence: what problem does a queue solve that a shared global
          variable does not?
  \end{enumerate}
\end{questionbox}

Restore \texttt{vDisplayTask} (remove the added \api{vTaskDelay}).

\subsection{Experiment 2 --- The mutex and the race condition}

\api{xUartMutex} ensures \texttt{Display} and \texttt{Status} never write the UART at
the same time.  Remove it and force a collision.

\begin{enumerate}[noitemsep, leftmargin=*]
  \item In \file{lab\_tasks.c}, comment out \textbf{all four}
        \api{xSemaphoreTake(xUartMutex, ...)} and \api{xSemaphoreGive(xUartMutex)}
        calls in \texttt{vDisplayTask} and \texttt{vStatusTask} (keep the
        \api{uart\_printf} lines).
  \item To make collisions frequent, speed both producers up:
        \begin{itemize}[noitemsep]
          \item \texttt{vSensorTask}: change \texttt{xPeriod} to
                \api{pdMS\_TO\_TICKS(50)}.
          \item \file{main.c}: change the timer period to \api{pdMS\_TO\_TICKS(100)}.
        \end{itemize}
\end{enumerate}

Rebuild, flash, and study the terminal.

\begin{questionbox}[~2]
  \begin{enumerate}[label=(\alph*), noitemsep, leftmargin=*]
    \item Do you see garbled lines --- a \texttt{[DISPLAY]} fragment spliced into the
          middle of a \texttt{[STATUS]} line, or vice versa?
    \item This is a \textbf{race condition}.  Explain why it happens \emph{even on a
          single-core MCU}.  (Hint: the UART sends one character at a time; the
          scheduler can preempt \texttt{Display} mid-string.)
    \item FreeRTOS mutexes implement \textbf{priority inheritance}
          (\texttt{configUSE\_MUTEXES~1}), but binary semaphores do not.  Why does that
          matter when a low-priority task holds a lock that a high-priority task wants?
  \end{enumerate}
\end{questionbox}

Restore the mutex calls and the original periods.

\subsection{Experiment 3 --- Binary vs.\ counting semaphore (lost signals)}

\api{xStatusSem} is a \textbf{binary} semaphore: it holds at most one ``token''.  If
the timer gives it again before \texttt{vStatusTask} takes it, the extra give is
\textbf{discarded}.

\begin{enumerate}[noitemsep, leftmargin=*]
  \item In \file{main.c}, speed up the timer: \api{pdMS\_TO\_TICKS(1000)} $\to$
        \api{pdMS\_TO\_TICKS(200)}.
  \item In \file{lab\_tasks.c}, slow down \texttt{vStatusTask} by adding a delay after
        the print:
\end{enumerate}

\begin{lstlisting}[style=cstyle]
xSemaphoreGive(xUartMutex);
vTaskDelay(pdMS_TO_TICKS(1000));   /* <-- ADD */
\end{lstlisting}

Rebuild and flash.  The timer now fires $\sim$5\,$\times$/s, but \texttt{vStatusTask}
can only process $\sim$1\,$\times$/s.

\begin{questionbox}[~3a]
  How many \texttt{[STATUS]} lines appear per second?  How many timer ``gives'' are
  being \textbf{lost}?
\end{questionbox}

Now switch to a \textbf{counting} semaphore.  In \file{main.c}:

\begin{lstlisting}[style=cstyle]
xStatusSem = xSemaphoreCreateCounting(10, 0);  /* was xSemaphoreCreateBinary() */
\end{lstlisting}

(\texttt{configUSE\_COUNTING\_SEMAPHORES} is already \texttt{1} in
\file{FreeRTOSConfig.h}.)  Rebuild and flash.

\begin{questionbox}[~3b]
  With the counting semaphore, what happens when \texttt{vStatusTask} finally runs
  after its 1000\,ms delay?  Why does it now print several lines back-to-back?  When is
  a binary semaphore the \emph{correct} choice, and when do you need a counting one?
\end{questionbox}

Restore the binary semaphore, the 1000\,ms timer period, and \texttt{vStatusTask}.

\subsection{Experiment 4 --- Software timer: one-shot vs.\ auto-reload}

The \texttt{StatusTimer} is created with \texttt{pdTRUE} --- \textbf{auto-reload}.
Change it to one-shot.  In \file{main.c}, in the \api{xTimerCreate} call:

\begin{lstlisting}[style=cstyle]
xTimerCreate("StatusTimer", pdMS_TO_TICKS(1000),
             pdFALSE,   /* <-- was pdTRUE: now one-shot */
             NULL, vStatusTimerCallback);
\end{lstlisting}

Rebuild and flash.

\begin{questionbox}[~4]
  \begin{enumerate}[label=(\alph*), noitemsep, leftmargin=*]
    \item How many \texttt{[STATUS]} lines do you see now, and why?
    \item The callback \api{vStatusTimerCallback} runs in the \textbf{timer service
          task} (\texttt{configTIMER\_TASK\_PRIORITY = 4}, the highest priority).
          Why must a timer callback \textbf{never} call a blocking API such as
          \api{vTaskDelay} or a \texttt{Take} with a long timeout?
    \item Name one way to make a one-shot timer fire repeatedly \emph{without} using
          \texttt{pdTRUE}.  (Hint: \api{xTimerStart} / \api{xTimerReset} can be called
          from inside the callback.)
  \end{enumerate}
\end{questionbox}

Restore \texttt{pdTRUE}.

\subsection{Challenge (optional) --- \texttt{vTaskDelay} vs.\ \texttt{vTaskDelayUntil}}

\texttt{vSensorTask} uses \api{vTaskDelayUntil} for a \textbf{drift-free} 500\,ms
period.  Replace it with \api{vTaskDelay} and add a variable workload:

\begin{lstlisting}[style=cstyle]
/* Simulate variable processing time before the delay */
for (volatile int i = 0; i < 200000; i++) { }

vTaskDelay(pdMS_TO_TICKS(500));        /* <-- was vTaskDelayUntil(...) */
\end{lstlisting}

Rebuild, flash, and look at the \texttt{t=} timestamps in the \texttt{[DISPLAY]} lines.

\begin{questionbox}[~Challenge]
  With \api{vTaskDelay}, the period becomes \emph{work time + 500\,ms}, so the
  timestamps drift away from clean multiples of 500.  With \api{vTaskDelayUntil} they
  stay locked to the 500\,ms grid.  Explain why --- what does \api{vTaskDelayUntil}
  measure its delay \emph{from} that \api{vTaskDelay} does not?
\end{questionbox}

Restore the original \api{vTaskDelayUntil} code.

% ── Deliverables ───────────────────────────────────────────────────────────────
\section{Deliverables}

Submit the following to the course LMS by the posted deadline:

\renewcommand{\arraystretch}{1.3}
\begin{tabular}{@{} c p{4.2cm} p{7.8cm} @{}}
  \toprule
  \textbf{\#} & \textbf{Item} & \textbf{Details} \\
  \midrule
  1 & Baseline capture     & 5+ seconds of correct output (Checkpoint~A) \\
  2 & Experiment 1 capture & Terminal output showing a \texttt{Queue full} drop \\
  3 & Experiment 2 capture & Terminal output showing garbled (interleaved) lines \\
  4 & Experiment 3 captures & Before \emph{and} after the counting semaphore \\
  5 & Answers to Q1--Q4    & A few sentences each; reference the code or output \\
  6 & Challenge answer     & Optional, if attempted \\
  \bottomrule
\end{tabular}
\renewcommand{\arraystretch}{1}

% ── Key Concepts ───────────────────────────────────────────────────────────────
\section{Key Concepts Introduced}

\renewcommand{\arraystretch}{1.3}
\begin{tabular}{@{} p{7cm} p{6.5cm} @{}}
  \toprule
  \textbf{Concept} & \textbf{Where you saw it} \\
  \midrule
  Producer--consumer decoupling via a queue & Experiment~1 \\
  Queue overflow and non-blocking \api{xQueueSend} & Experiment~1 \\
  Race condition on a shared resource & Experiment~2 \\
  Mutex as a critical-section guard; priority inheritance & Experiment~2 \\
  Binary vs.\ counting semaphore; lost signals & Experiment~3 \\
  Software timers: one-shot vs.\ auto-reload; callback rules & Experiment~4 \\
  Drift-free periodic timing & Challenge \\
  \bottomrule
\end{tabular}
\renewcommand{\arraystretch}{1}

\medskip
These primitives are the toolkit for Module~3 (synchronisation) and Module~4
(inter-task communication).

% ── Reference ─────────────────────────────────────────────────────────────────
\section{Reference}

\begin{tabular}{@{} p{6cm} p{7.5cm} @{}}
  \toprule
  \textbf{Resource} & \textbf{Relevant sections} \\
  \midrule
  Barry, \textit{Mastering the FreeRTOS Kernel} &
    Ch.~4 (queues), Ch.~6 (semaphores/mutexes), Ch.~5 (software timers) \\
  Laplante, \textit{Real-Time Systems Design} &
    Ch.~3 (intertask communication, synchronisation) \\
  FreeRTOS documentation &
    \api{xQueueCreate}, \api{xQueueSend}, \api{xQueueReceive},
    \api{xSemaphoreCreateMutex}, \api{xSemaphoreCreateBinary},
    \api{xSemaphoreCreateCounting}, \api{xTimerCreate} \\
  \bottomrule
\end{tabular}

\end{document}
